% =====================================================================
\chapter{Входы: расходы, экспонента, триммеры}
% =====================================================================
\chapterintro
  {как стик превращается во «вход», что такое расходы (dual rates) и экспонента, как сделать
   переключаемые расходы и зачем нужна выборка по полётным режимам}
  {модель с настроенным радиомодулем}

\section{Страница Inputs}

\mpath{MDL\sep Inputs} содержит список входов \sw{I1}, \sw{I2}, … У новой модели EdgeTX
создаёт четыре входа по количеству осей стиков, в порядке \menu{Default channel order}:

\begin{center}
\lcd{\inv{INPUTS}\hfill 4/64\\
{[I1]}Ail\quad Ail\quad 100\\
{[I2]}Ele\quad Ele\quad 100\\
{[I3]}Thr\quad Thr\quad 100\\
{[I4]}Rud\quad Rud\quad 100}
\end{center}

Каждый вход может состоять из \emph{нескольких строк}. EdgeTX проверяет строки сверху вниз
и использует \textbf{первую}, у которой выполнено условие (переключатель, полётный режим).
На этом построены переключаемые расходы.

\section{Параметры строки входа}

Нажмите \btn{ENTER} на строке → \menu{Edit}:

\begin{center}\small
\renewcommand{\arraystretch}{1.2}
\begin{tabularx}{\linewidth}{@{}p{30mm}X@{}}
\toprule
\menu{Input name} & имя входа (одно на весь вход) \\
\menu{Line name} & имя конкретной строки, например \code{High}, \code{Low} \\
\menu{Source} & источник: стик, крутилка, переключатель, канал, телеметрия \\
\menu{Weight} & вес --- масштаб, $-100\ldots+100\,\%$. Это и есть «расход» (rate) \\
\menu{Offset} & смещение, прибавляется после умножения на вес \\
\menu{Trim} & учитывать ли триммер источника (\menu{ON}), не учитывать (\menu{OFF}) или
  взять триммер другой оси \\
\menu{Curve} & функция формы: \menu{Diff} (дифференциал), \menu{Expo} (экспонента),
  \menu{Func} (x>0, |x|, …), \menu{Curve} (своя кривая) \\
\menu{Flight modes} & в каких полётных режимах строка активна \\
\menu{Switch} & условие активности строки \\
\menu{Side} & работать на всём ходе стика, только на $x>0$ или только на $x<0$ \\
\bottomrule
\end{tabularx}
\end{center}

\section{Расходы (dual rates)}

\textbf{Расход} --- максимальный ход, который даёт стик. Вес 100\,\% --- полный ход,
60\,\% --- стик до упора даёт только 60\,\% отклонения рулей. Малые расходы используют
для первых полётов и точного пилотажа, большие --- для фигур.

\section{Экспонента}

\textbf{Экспонента (expo)} делает управление мягче у центра стика, сохраняя полный ход
по краям. В EdgeTX формула такая (для нормированного хода $x \in [-1, 1]$ и значения
экспоненты $k = \text{expo}/100$):
\[
  y = x\,\bigl(k x^2 + (1-k)\bigr).
\]

\begin{figure}[H]
\centering
\begin{tikzpicture}
\begin{axis}[width=9cm,height=6.5cm,xmin=-100,xmax=100,ymin=-100,ymax=100,
  xlabel={ход стика, \%},ylabel={выход, \%},grid=major,grid style={gray!20},
  legend pos=north west,legend style={font=\scriptsize},tick label style={font=\scriptsize},
  label style={font=\small}]
\addplot[domain=-100:100,samples=60,thick,gray] {x};
\addplot[domain=-100:100,samples=60,thick,etx] {x*(0.3*(x/100)^2+0.7)};
\addplot[domain=-100:100,samples=60,thick,accent] {x*(0.6*(x/100)^2+0.4)};
\addplot[domain=-100:100,samples=60,thick,okgreen,dashed] {0.7*x*(0.3*(x/100)^2+0.7)};
\legend{expo 0,expo 30\,\%,expo 60\,\%,вес 70\,\%\,+\,expo 30\,\%}
\end{axis}
\end{tikzpicture}
\caption{Экспонента и расход: у центра стика управление мягче, края --- без изменений}
\end{figure}

Типичные значения: 20--40\,\% на элеронах и руле высоты самолёта. Для квадрокоптера с
Betaflight экспоненту и расходы обычно задают в полётном контроллере (rates),
а в пульте оставляют вес 100\,\% и экспоненту 0.

\section{Переключаемые расходы}

\begin{recipe}{три расхода на элеронах по переключателю SB}
Вход \sw{I1} (Ail) делаем из трёх строк:
\begin{center}\small
\begin{tabular}{@{}llllll@{}}
\toprule
Строка & Source & Weight & Curve & Switch \\
\midrule
High   & Ail & 100 & Expo 25 & \sw{SB↑} \\
Mid    & Ail & 75  & Expo 30 & \sw{SB-} \\
Low    & Ail & 50  & Expo 35 & --- (всегда) \\
\bottomrule
\end{tabular}
\end{center}
Порядок важен: последняя строка без переключателя срабатывает, если ни одна верхняя не
подошла. Точно так же настройте \sw{I2} (Ele) на том же переключателе.
\end{recipe}

\begin{center}
\begin{tikzpicture}[node distance=3mm,every node/.style={font=\sffamily\scriptsize}]
  \node[blk,fill=blksw] (stk) {стик Ail};
  \node[blk,fill=blkin,right=12mm of stk,yshift=9mm] (h) {High: 100\,\%};
  \node[blk,fill=blkin,right=12mm of stk] (m) {Mid: 75\,\%};
  \node[blk,fill=blkin,right=12mm of stk,yshift=-9mm] (l) {Low: 50\,\%};
  \node[blk,fill=blkmix,right=14mm of m] (i1) {\sw{I1}};
  \draw[arr] (stk.east) -- (h.west);
  \draw[arr] (stk.east) -- (m.west);
  \draw[arr] (stk.east) -- (l.west);
  \draw[arr] (h.east) -- node[above,sloped,lbl]{SB↑} (i1.west);
  \draw[arr] (m.east) -- node[above,lbl]{SB-} (i1.west);
  \draw[arr] (l.east) -- node[below,sloped,lbl]{иначе} (i1.west);
\end{tikzpicture}
\end{center}

\section{Дифференциал и функции}

\begin{itemize}
  \item \menu{Diff} уменьшает ход в одну сторону. На входе используется редко; дифференциал
        элеронов правильнее задавать в миксерах (глава~\ref{ch:mixes}).
  \item \menu{Func}: \menu{x>0} --- пропускать только положительную половину хода,
        \menu{|x|} --- модуль, \menu{f>0}, \menu{sgn} (знак) и др. Например, для кнопки
        аварийного сброса или для «половинного» стика.
\end{itemize}

\section{Триммеры во входах}

Параметр \menu{Trim} = \menu{ON} (по умолчанию) добавляет значение триммера оси.
Для квадрокоптеров выключайте триммеры (\menu{OFF}) --- случайное нажатие качельки изменит
центр стика, и полётный контроллер будет считать, что вы держите стик отклонённым.

\begin{tasks}
\begin{enumerate}
  \item Сделайте для самолёта два расхода на \sw{SC}: 100\,\% с экспо 20\,\% и 60\,\% с экспо 30\,\%.
  \item Проверьте работу на странице \menu{Outputs} или на экране каналов: переключайте \sw{SC}
        и смотрите, как меняется максимум.
  \item Сделайте вход \sw{I5} с источником \sw{S1} и весом 50\,\% --- «половинную» крутилку.
\end{enumerate}
\end{tasks}
